\section{Diffusion de Perona--Malik : adaptation et difficultés}
\label{sec:perona-malik}

\subsection{Un coefficient dépendant des variations locales}

Le modèle continu proposé par Perona et Malik est
\begin{equation}
 \partial_t u=\operatorname{div}\bigl(c(|\nabla u|)\nabla u\bigr),
 \qquad u(\cdot,0)=f,
 \qquad c(|\nabla u|)\nabla u\cdot\boldsymbol n=0.
 \label{eq:pm-continuous}
\end{equation}
Il vise à favoriser le lissage des régions peu variables tout en
réduisant les échanges au voisinage des fortes variations
\cite{perona1990}. Les deux lois étudiées sont
\begin{equation}
 c_{\mathrm{exp}}(s)=\exp\left(-\frac{s^2}{K^2}\right),\qquad
 c_{\mathrm{rat}}(s)=\frac{1}{1+s^2/K^2},\qquad s\geq0,
 \quad K>0.
 \label{eq:pm-conductances}
\end{equation}
Elles valent un en $s=0$, sont décroissantes pour $s>0$ et restent
strictement positives pour une valeur finie de $s$ dans le modèle
mathématique. Dans le calcul flottant, l'exponentielle peut devenir
numériquement nulle pour un très grand rapport $s/K$ ; cette valeur
respecte encore la borne $0\leq c\leq1$ utilisée dans l'analyse discrète.

Le paramètre $K$ est une échelle de variation d'intensité par unité de
longueur. Avec un pas spatial égal à un pixel, il se lit dans la
convention intensité par pixel. Ce n'est ni une probabilité de contour
ni un seuil de segmentation. Pour $s\ll K$, les deux coefficients sont
proches de un et le comportement local se rapproche de la chaleur.
Pour $s\gg K$, la diffusion est freinée. L'exponentielle décroît plus
rapidement que la loi rationnelle ; à $s=K$, les valeurs sont
respectivement $\mathrm e^{-1}$ et $1/2$.

La distinction entre contour et bruit n'est pas donnée au modèle.
Une différence importante produite par le bruit peut également
réduire le coefficient de diffusion. Un petit $K$ peut ainsi conserver
des fluctuations que l'on souhaite éliminer. À l'inverse, un grand $K$
rend davantage de transitions perméables et accroît le risque de
lissage des structures. Comme les coefficients sont recalculés à
chaque étape, leur effet dépend aussi du temps : les variations qui
diminuent peuvent ensuite autoriser plus de diffusion. La comparaison
doit donc porter conjointement sur $K$, la loi $c$ et l'arrêt, sans
présupposer la supériorité d'une conductance sur l'autre.

\subsection{Pourquoi une conductance positive ne suffit pas}

Posons $g(s)=s c(s)$, amplitude du flux dans la direction du gradient.
Une conductance positive n'implique pas que $g$ soit croissant.
Le calcul direct donne
\begin{align}
 g'_{\mathrm{exp}}(s)
 &=\exp(-s^2/K^2)\left(1-2s^2/K^2\right),
 \label{eq:pm-flux-exp}\\
 g'_{\mathrm{rat}}(s)
 &=\frac{1-s^2/K^2}{(1+s^2/K^2)^2}.
 \label{eq:pm-flux-rat}
\end{align}
La dérivée est nulle en $s=K/\sqrt2$ pour l'exponentielle et en $s=K$
pour la rationnelle ; elle est négative au-delà. Ces seuils concernent
la monotonie du flux, non la valeur à partir de laquelle la conductance
s'annule. Les deux lois possèdent donc un régime où une augmentation
de la pente diminue l'amplitude du flux.

Pour préciser la difficulté en dimension deux, écrivons
$A(p)=c(|p|)p$, $p\in\R^2$. Pour $p\ne0$ et $s=|p|$, sa matrice
différentielle est
\begin{equation}
 D A(p)=c(s)I+\frac{c'(s)}{s}p p^{\mathsf T}.
 \label{eq:pm-jacobian}
\end{equation}
Un vecteur orthogonal à $p$ est vecteur propre de valeur $c(s)$,
tandis que $p$ est vecteur propre de valeur $c(s)+s c'(s)=g'(s)$.
Dans une linéarisation autour d'un gradient suffisamment fort,
la seconde valeur devient négative. L'opérateur n'est donc pas
uniformément parabolique dans tous les régimes. Le problème continu
présente des difficultés de type diffusion avant--arrière ; un
argument fondé seulement sur $c>0$ serait insuffisant.

Ces dérivées sont ici démontrées par calcul élémentaire. Elles
éclairent la différence entre la motivation de préservation des
structures et les exigences d'une théorie de bonne position. Le fait
que les images numériques demeurent bornées sur une grille finie ne
prouve ni existence ni unicité ni dépendance continue pour toutes
les données du problème continu \eqref{eq:pm-continuous}.

\subsection{Une interprétation variationnelle formelle}

Une primitive adaptée est définie par $\Psi'(s)=s c(s)$ et
$\Psi(0)=0$. L'intégration donne
\begin{equation}
 \Psi_{\mathrm{exp}}(s)=\frac{K^2}{2}
    \left(1-\exp(-s^2/K^2)\right),\qquad
 \Psi_{\mathrm{rat}}(s)=\frac{K^2}{2}
    \log\left(1+s^2/K^2\right).
 \label{eq:pm-potentials}
\end{equation}
Pour une fonction suffisamment régulière et des variations compatibles
avec le flux nul, le premier accroissement de
\begin{equation}
 E_\Psi(u)=\int_\Omega\Psi(|\nabla u|)\,\dd x\,\dd y
 \label{eq:pm-energy-formal}
\end{equation}
est
\begin{equation}
 \dd E_\Psi(u)[v]
 =\int_\Omega c(|\nabla u|)\nabla u\cdot\nabla v
 =-\int_\Omega
    \operatorname{div}\bigl(c(|\nabla u|)\nabla u\bigr)v.
 \label{eq:pm-firstvariation}
\end{equation}
Le modèle est donc formellement un flot de gradient négatif pour
cette fonctionnelle. Le calcul s'étend au point $\nabla u=0$ puisque
$\Psi'(0)=0$ et les deux densités sont régulières comme fonctions
du vecteur gradient.

Cependant, la Hessienne de $p\mapsto\Psi(|p|)$ est précisément
$D A(p)$. Ses valeurs propres sont $c(s)$ dans les directions
tangentielles et $g'(s)$ dans la direction radiale. D'après
\eqref{eq:pm-flux-exp}--\eqref{eq:pm-flux-rat}, ces densités ne sont
pas convexes globalement. Elles ne fournissent pas le cadre convexe
de l'énergie de Dirichlet. Une identité de dissipation formelle, si
une solution classique et les opérations de dérivation existent,
ne construit pas cette solution et ne résout pas les difficultés
de passage à la limite. Elle ne justifie pas non plus une décroissance
de l'énergie pour l'algorithme explicite réellement utilisé, dont la
conductance est directionnelle. Aucune telle décroissance non linéaire
n'est revendiquée dans ce projet.

\subsection{Ce qui est implémenté et ce qui reste une perspective}

Le modèle \eqref{eq:pm-continuous} utilise la norme complète du
gradient dans un coefficient scalaire. Le solveur étudié emploie,
sur chaque interface, uniquement la différence dans la direction
normale à cette interface. Il s'agit de la \emph{variante discrète
directionnelle à quatre voisins}, détaillée dans la section suivante.
Cette dénomination doit accompagner l'interprétation des résultats.
Les motivations du modèle original restent pertinentes, mais le
schéma ne doit pas être présenté comme une reconstruction isotrope
exacte de $c(|\nabla u|)$.

Catté, Lions, Morel et Coll proposent de régulariser la dépendance
du coefficient en évaluant les variations sur une image spatialement
lissée \cite{catte1992}. Cette approche sépare l'image diffusée du
signal utilisé pour guider sa diffusion. Elle constitue une
perspective bibliographique pour traiter les difficultés du modèle,
et non une méthode expérimentée ici. Aucun noyau de régularisation
ni paramètre de lissage supplémentaire n'a été ajouté aux deux
solveurs étudiés.
